\section{Introduction}
\label{sec:intro}

Machine-learning network intrusion detection systems (NIDS) classify network flows as benign or as one of several
attack classes. On public benchmarks such as CIC-IDS2017~\cite{sharafaldin2018} they routinely report accuracies
above 99\%. Most of these results come from random train/test splits, in which flows of the same attack session
fall on both sides of the split, so the test asks whether the model recognises traffic it has effectively already
seen. Network traffic is time-dependent, and estimates obtained with a chronological split differ from those of a
random split~\cite{nismara2026temporal}. Careful sequence evaluations find that the evaluation protocol matters far
more than the model architecture~\cite{moczkodan2026transformers}. These are instances of the temporal and spatial
biases described for security classifiers in general~\cite{pendlebury2019tesseract,arp2022dos}.
\Cref{sec:res-split} quantifies the gap on our data. The same gradient-boosting classifier reaches 99.8\%
accuracy on a random split of CIC-IDS2017, the figure usually reported. Trained on the first 30\% of the capture
and tested on the rest, it reaches 72.6\%, and its macro-F1 falls from 0.761 to 0.125, because eleven of the
classes in the test period never occur in training.

A deployed detector faces four conditions that random splits hide. First, \emph{new attack classes} appear after
training. A closed-set model assigns them to some known class, most often to benign traffic, and nothing tells
the operator that it is guessing. Second, \emph{labels are scarce}: an analyst can inspect only a small fraction of
the flows. Third, \emph{labels are late}: ground truth for today's traffic arrives after today's decisions were
made. Fourth, every model update is a \emph{deployment}, and a bad update degrades detection silently unless
something checks it. Prior work addresses these conditions one at a time (\cref{sec:related}). Drift
detectors~\cite{jordaney2017transcend,yang2021cade} flag drifting samples but do not adapt. Adaptive
NIDS~\cite{andresini2021insomnia,zhang2024aoc,zhang2025ssf} update mostly binary detectors, often from their own
pseudo-labels, and rarely model label delay. Active-learning methods~\cite{chen2023continuous,gupta2025netguard}
choose which samples to label but assume a closed set of classes. Open-set NIDS~\cite{zhang2020osr,corsini2023ood}
reject unknown traffic, but only on static splits.

This article studies the four conditions together and asks five research questions:
\begin{description}
  \item[RQ1] How large is the gap between random-split and chronological evaluation for the same classifier?
  \item[RQ2] Under one label budget and delay, which adaptation strategy keeps a multi-class open-world detector
        most accurate, and at what false-positive cost?
  \item[RQ3] How do the strategies behave when labels become ten times scarcer?
  \item[RQ4] Which components of the proposed method matter, and do their effects carry over to datasets that were
        not used for its design?
  \item[RQ5] How sensitive are the conclusions to the label delay, the rejection thresholds and the classifier?
\end{description}

\textbf{Contributions.}
\begin{enumerate}
  \item An evaluation protocol for drift-aware open-world NIDS: chronological replay in windows, a label oracle
        that enforces the same budget and delay for every method, and metrics that stay comparable when methods
        learn different classes (\cref{sec:problem}). We show a property of open-world scoring that matters in
        practice: a method that never learns a new class can outscore one that learns it imperfectly
        (\cref{sec:res-budget}).
  \item \nova{} (Novelty-aware, Open-world, Validated Adaptation), which combines class-conditional $k$-NN
        rejection with per-class recalibration, a novelty-weighted and cluster-diverse query, agreement-checked
        label propagation, a benign-weighted class-balanced memory, and a deployment guard with rollback
        (\cref{sec:method}).
  \item A comparison with eight methods, including re-implementations of four published approaches, on three
        datasets from two laboratories, with all design decisions made on one development dataset and the other
        two used only for evaluation (\cref{sec:setup,sec:results}).
  \item Analyses that explain the results: a full ablation on all three datasets, per-class results, sensitivity
        to label delay, rejection thresholds and classifier, and an offline test of whether a meta-learned
        representation would help (\cref{sec:results}).
  \item An open-source research console that stores the configuration, data hash and per-flow decision of every
        run, so that each number in this article can be reproduced.
\end{enumerate}

\textbf{Relation to our conference paper.} \TBD{A preliminary version of this work appears in [conference
reference]. This article adds the random-versus-chronological comparison (RQ1), a dataset audit, per-window
confidence intervals and per-class explanations for the ablations, the metric decomposition and guard statistics,
the label-delay, threshold and classifier sensitivity studies, per-class and failure analyses, the
representation-learning probe, an extended related-work survey and a full description of the protocol and
method.}
